5.7 Effect of the Risk-Aware Reward

Earlier experiments showed that the main wealth-change reward did not
always encourage agents to respond differently to changes in the
observed state. It simply rewards changes in portfolio wealth.

Chapter 3.5 identifies the limitation of this objective: it does not
penalize large portfolio drawdowns, raising the question of whether
explicitly including portfolio drawdown in the reward would lead to more
risk-aware behaviour.

To examine this, the risk-aware reward was introduced in which it was
modified to include a penalty for increases in drawdown:

\[r_{t}^{risk} = \Delta W_{t} - \lambda C_{0}\Delta dd_{t}.\]

where

\(\Delta dd_{t}\) is the increase in portfolio drawdown at trading day
\(t\) (and is zero when drawdown does not rise), and coefficient
\(\lambda\) controls the penalty effects. The initial capital \(C_{0}\)
factor puts \(\Delta dd_{t}\) in dollars so that 𝜆 is dimensionless.

The risk-aware reward is paired with the ten-feature state using the
drawdown feature because it provides the agent with information about
its current position relative to its previous maximum wealth.

\textbf{5.7.1 Selection of the Penalty size}

The penalty size was selected using the validation set and Deep
Q-learning as the selector because it is the only algorithm whose
policies consistently respond to the state (Section 5.4). The held-out
test period was not used for this decision.

A selection rule was created. The penalised configuration must keep at
least half (50\%) of both the unpenalised return---to still earn a
decent amount --- and the unpenalised exposure---to still participate in
the market.

The unpenalised DQN configuration had a mean wealth change of \$78.35
and exposure of 0.445. The resulting minimum requirements were
therefore:

\[\Delta W \geq \$ 39.17\]

and

\[Exposure \geq 0.223.\]

The results of the penalty sweep are shown in Table 5.5.

{\def\LTcaptype{none} % do not increment counter
\begin{longtable}[]{@{}
  >{\raggedright\arraybackslash}p{(\linewidth - 10\tabcolsep) * \real{0.1900}}
  >{\centering\arraybackslash}p{(\linewidth - 10\tabcolsep) * \real{0.1610}}
  >{\centering\arraybackslash}p{(\linewidth - 10\tabcolsep) * \real{0.1633}}
  >{\centering\arraybackslash}p{(\linewidth - 10\tabcolsep) * \real{0.1652}}
  >{\centering\arraybackslash}p{(\linewidth - 10\tabcolsep) * \real{0.1605}}
  >{\centering\arraybackslash}p{(\linewidth - 10\tabcolsep) * \real{0.1600}}@{}}
\toprule\noalign{}
\begin{minipage}[b]{\linewidth}\centering
\textbf{λ}
\end{minipage} & \begin{minipage}[b]{\linewidth}\centering
\textbf{Mean ΔW}
\end{minipage} & \begin{minipage}[b]{\linewidth}\centering
\textbf{Exposure}
\end{minipage} & \begin{minipage}[b]{\linewidth}\centering
\textbf{Drawdown}
\end{minipage} & \begin{minipage}[b]{\linewidth}\centering
\textbf{DD change}
\end{minipage} & \begin{minipage}[b]{\linewidth}\centering
\textbf{Passes}
\end{minipage} \\
\midrule\noalign{}
\endhead
\bottomrule\noalign{}
\endlastfoot
\textbf{0 (unpenalised)} & +\$78.35 & 0.445 & 0.0163 & --- & --- \\
\textbf{0.05} & +\$88.78 & 0.401 & 0.0128 & −21.5\% & yes \\
\textbf{0.10} & +\$67.30 & 0.344 & 0.0124 & −24.0\% & yes \\
\textbf{0.25} & +\$19.90 & 0.074 & 0.0035 & −78.4\% & no \\
\textbf{0.50} & +\$1.73 & 0.002 & 0.0001 & −99.4\% & no \\
\textbf{1.00} & \$0.00 & 0.000 & 0.0000 & −100\% & no \\
\textbf{2.00} & \$0.00 & 0.000 & 0.0000 & −100\% & no \\
\end{longtable}
}

\emph{Table 5.5: Risk-penalty size on the validation months (Deep
Q-learning, six seeds). The rule requires ΔW ≥ +\$39.17 and exposure ≥
0.223; the largest passing value that meets both requirements is λ =
0.10.}

\includegraphics[width=5.9in,height=2.52857in]{media5/media/image1.png}

\emph{Figure 5.6: Results from the penalty size of Table 5.5. Left: mean
wealth change against λ with the requirement floor of +\$39.17. Right:
exposure and drawdown relative to their values without penalties.
Between λ = 0.10 and 0.25, exposure drops sharply rather than decreasing
gradually.}

The results show a clear change in behaviour of the penalty \(\lambda\)
as the size increases. At \(\lambda = 0.10\) , the agent still maintains
a meaningful level of exposure and a return of \$67.30.

Increasing the penalty to 0.25 causes exposure to fall sharply to 0.074.
The agent loses about 78\% of its exposure, indicating reluctance to
trade. At \(\lambda = 0.50\) and above, the agent stops trading
completely.

This result depicts an important limitation of the reward formulation.
The simplest way for an agent to reduce drawdown is to hold less of the
asset or simply not hold it at all. In theory, a portfolio that remains
in cash is not exposed to asset movements and therefore does not
experience market drawdowns. Once the drawdown penalty becomes too
large, market participation becomes unattractive.

Based on these criteria, the largest coefficient that satisfied both the
return and exposure requirements was therefore:

\[\boxed{\lambda = 0.10}.
\]This became the fixed penalty value for evaluating the main risk-aware
experiment.

\textbf{5.7.2 Effect on the Test Set}

The selected risk-aware reward was then evaluated on the held-out test
set.

{\def\LTcaptype{none} % do not increment counter
\begin{longtable}[]{@{}
  >{\raggedright\arraybackslash}p{(\linewidth - 14\tabcolsep) * \real{0.1675}}
  >{\centering\arraybackslash}p{(\linewidth - 14\tabcolsep) * \real{0.1192}}
  >{\centering\arraybackslash}p{(\linewidth - 14\tabcolsep) * \real{0.1202}}
  >{\centering\arraybackslash}p{(\linewidth - 14\tabcolsep) * \real{0.1162}}
  >{\centering\arraybackslash}p{(\linewidth - 14\tabcolsep) * \real{0.1364}}
  >{\centering\arraybackslash}p{(\linewidth - 14\tabcolsep) * \real{0.1125}}
  >{\centering\arraybackslash}p{(\linewidth - 14\tabcolsep) * \real{0.1117}}
  >{\centering\arraybackslash}p{(\linewidth - 14\tabcolsep) * \real{0.1162}}@{}}
\toprule\noalign{}
\begin{minipage}[b]{\linewidth}\centering
\textbf{Algorithm}
\end{minipage} & \begin{minipage}[b]{\linewidth}\centering
\textbf{Reward}
\end{minipage} & \begin{minipage}[b]{\linewidth}\centering
\textbf{ΔW (\$)}
\end{minipage} & \begin{minipage}[b]{\linewidth}\centering
\textbf{Spread}
\end{minipage} & \begin{minipage}[b]{\linewidth}\centering
\textbf{Exposure}
\end{minipage} & \begin{minipage}[b]{\linewidth}\centering
\textbf{Trades}
\end{minipage} & \begin{minipage}[b]{\linewidth}\centering
\textbf{DD}
\end{minipage} & \begin{minipage}[b]{\linewidth}\centering
\textbf{Sharpe}
\end{minipage} \\
\midrule\noalign{}
\endhead
\bottomrule\noalign{}
\endlastfoot
\textbf{REINFORCE} & wealth & +167.39 & 24.7 & 0.794 & 13.0 & 0.0263 &
0.670 \\
\textbf{REINFORCE} & risk & +167.39 & 24.7 & 0.794 & 13.0 & 0.0263 &
0.670 \\
\textbf{Deep Q} & wealth & +77.81 & 89.5 & 0.354 & 7.4 & 0.0140 &
0.505 \\
\textbf{Deep Q} & risk & +52.90 & 43.1 & 0.178 & 5.4 & 0.0082 & 0.416 \\
\textbf{PPO} & wealth & +159.97 & 94.0 & 0.759 & 7.3 & 0.0250 & 0.671 \\
\textbf{PPO} & risk & +21.48 & 128.9 & 0.108 & 0.7 & 0.0037 & 0.107 \\
\end{longtable}
}

\emph{Table 5.6: Wealth-change versus risk-aware reward (λ = 0.10) on
the held-out months, ten-feature state, six seeds.}

The DQN result gives the clearest evidence of the effect of the
risk-aware reward. For Deep Q-learning, the mean drawdown falls from
0.0140 under the standard reward to 0.0082 with the risk-aware reward, a
reduction of approximately 41\%. However, its return also falls from
+\$77.81 to +\$52.90, a reduction of about 32\%. Exposure falls from
0.354 to 0.178, meaning that the agent decides to hold only half as much
of the asset.

The results do not suggest improved risk-adjusted performance. Instead,
the reward mainly encouraged the agent to reduce its exposure to the
asset, and the Sharpe ratio also decreased from 0.505 to 0.416.

The PPO results show the same effect more strongly. With the risk-aware
reward, PPO\textquotesingle s exposure fell to 0.108, and its average
trading activity dropped below one trade per month, from 7.3 to 0.7. As
a result, its mean wealth change decreased from \$159.97 to +\$21.48.

REINFORCE shows no noticeable change. Its results are identical for both
rewards. This is consistent with the earlier behavioural findings:
REINFORCE policies do not change their trading behaviour simply by
changing the reward.

\textbf{5.7.3 Interpretation of the risk-aware return}

The risk-aware reward changed the behaviour of the state-dependent
agents, but there is no clear evidence that it led to better risk
management. Instead, the results suggest that the penalty mainly
encouraged the agents to reduce their exposure to the asset.

These results do not show that risk-aware reinforcement learning is
ineffective in general. However, they do show a limitation of the reward
formulation used in this study. An agent's policy could reduce drawdown
to zero by simply holding less of that asset or by never entering the
market, but this would not satisfy the main objective of generating
useful investment returns.

A more effective objective would therefore need to consider both return
and risk, rather than mainly penalising exposure to drawdown. The
relevant question becomes whether the agent can \textbf{reduce drawdown
while still leading to meaningful returns and exposure.} For example,
the objective could reward the agent for achieving higher returns for a
given level of risk. Chapter 6 discusses other possible alternatives in
future work.

\textbf{5.8 Transaction-Fee Sizing}

Transaction costs are an important part of the trading environment
because frequent trading can reduce the returns produced by a strategy,
even when individual decisions appear reasonable. A strategy may appear
profitable before fees are included but become less attractive once the
cost of each trade is taken into account.

The main experiments use a transaction fee of 0.05\%. To test whether
the observed behaviour changes as this fee varies, the fee was swept
from 0 basis points (0.00\%) to 50basis points (0.50\%) on the
validation months with the nine-feature state, and everything else held
fixed.

The sweep was run on all algorithms, each with six freshly trained seeds
per fee level. Table 5.7 and Figure 5.7 report the outcome.

{\def\LTcaptype{none} % do not increment counter
\begin{longtable}[]{@{}
  >{\raggedright\arraybackslash}p{(\linewidth - 14\tabcolsep) * \real{0.1320}}
  >{\centering\arraybackslash}p{(\linewidth - 14\tabcolsep) * \real{0.1243}}
  >{\centering\arraybackslash}p{(\linewidth - 14\tabcolsep) * \real{0.1242}}
  >{\centering\arraybackslash}p{(\linewidth - 14\tabcolsep) * \real{0.1237}}
  >{\centering\arraybackslash}p{(\linewidth - 14\tabcolsep) * \real{0.1242}}
  >{\centering\arraybackslash}p{(\linewidth - 14\tabcolsep) * \real{0.1237}}
  >{\centering\arraybackslash}p{(\linewidth - 14\tabcolsep) * \real{0.1242}}
  >{\centering\arraybackslash}p{(\linewidth - 14\tabcolsep) * \real{0.1237}}@{}}
\toprule\noalign{}
\begin{minipage}[b]{\linewidth}\centering
\textbf{Fee}
\end{minipage} & \begin{minipage}[b]{\linewidth}\centering
\textbf{B\&H ΔW}
\end{minipage} & \begin{minipage}[b]{\linewidth}\centering
\textbf{REINF. ΔW}
\end{minipage} & \begin{minipage}[b]{\linewidth}\centering
\textbf{REINF. trades}
\end{minipage} & \begin{minipage}[b]{\linewidth}\centering
\textbf{Deep Q ΔW}
\end{minipage} & \begin{minipage}[b]{\linewidth}\centering
\textbf{Deep Q trades}
\end{minipage} & \begin{minipage}[b]{\linewidth}\centering
\textbf{PPO ΔW}
\end{minipage} & \begin{minipage}[b]{\linewidth}\centering
\textbf{PPO trades}
\end{minipage} \\
\midrule\noalign{}
\endhead
\bottomrule\noalign{}
\endlastfoot
\textbf{0bps(0.00\%)} & +\$177.09 & +\$177.56 & 5.0 & +\$108.00 & 10.9 &
+\$155.35 & 4.5 \\
\textbf{5bps(0.05\%)} & +\$166.92 & +\$154.96 & 10.3 & +\$86.16 & 8.5 &
+\$118.55 & 3.7 \\
\textbf{10bps(0.10\%)} & +\$156.76 & +\$143.74 & 4.7 & +\$77.60 & 8.9 &
+\$78.50 & 2.7 \\
\textbf{25bps(0.25\%)} & +\$126.33 & +\$88.76 & 10.3 & +\$32.35 & 3.9 &
+\$26.40 & 1.2 \\
\textbf{50bps(0.50\%)} & +\$75.83 & +\$41.31 & 2.8 & +\$16.32 & 1.9 &
+\$12.60 & 1.0 \\
\end{longtable}
}

\emph{Table 5.7: Transaction Fee sizing on the validation months; six
seeds. Trades are per month. Fees are shown in basis points and as a
percentage of trade value.}

\includegraphics[width=6in,height=2.58462in]{media5/media/image2.png}

\emph{Figure 5.7: Transaction Fee sizing from Table 5.7. Left: average
wealth change falls for every strategy as fees rise. Right: average
trades per month as fees rise.}

\textbf{5.8.1 Effect of transaction costs on returns \& trading
behaviour}

As expected, higher transaction fees reduce the returns of the different
algorithm strategies. The effect is important for strategies that trade
frequently because more of their portfolio gains go to transaction
costs.

Deep Q-learning shows the clearest reduction in trading activity as the
fee of taking a trade increases. Its average number of trades falls from
10.9 per month at 0bps (0.00\%) to 1.9 at 50bps (0.50\%), an 83\%
reduction. As each action gets more expensive, the agent acts less,
reducing exposure. This is consistent with the algorithm correctly
incorporating transaction costs into its environment.

REINFORCE does not show the same clear pattern. Its trade counts vary,
running from 5.0, 10.3, 4.7, 10.3, and 2.8, respectively, with no
consistent decrease in trades taken as the fee rises. This aligns with
the earlier state-independent results: an agent that does not rely on
information from its current state to make decisions is less likely to
adjust its behaviour when trading costs change.

PPO's average trade count also falls from 4.5 to 1.0, showing that many
seeds are pushed to stop trading due to higher fees. This is also why
the average return collapses by 77\%, from +\$177.56 to +\$41.31, a
difference of \$136.25 per month, even though the average trade count
moves only gently.

\textbf{5.8.2 Interpretation}

The transaction-fee experiment therefore reflects another difference
between the learning algorithms in terms of their change in behaviour.
Deep Q-learning changes its trading behaviour as trading costs increase,
even though its overall returns are below the buy-and-hold benchmark.
REINFORCE produces returns generally closer to the buy-and-hold
strategy, and it also shows higher exposure that does not respond
consistently to changes in transaction costs. PPO, on the other hand,
shows a major collapse in returns as fees increase, while its trade
counts decline more gently. Overall, the learned agents lose
proportionally more because they trade more often in this experiment.

This indicates that the learning system is capable of responding to
different changes when the learned policy is sufficiently
state-dependent.

\textbf{5.9 Discussion of the Main Findings}

The experiments produced several findings about both the performance of
the learning agents and the way they behaved within the trading
environment. They also answer the research questions of Section 5.1

\textbf{5.9.1 The Benchmark Remains Difficult to Beat}

The main finding is that none of the learned agents could outperform the
simple buy-and-hold benchmark on the held-out test period. Buy-and-hold
achieved a wealth change of \$180.25, compared with \$171.51 for
REINFORCE as the best learned agent, while every other agent fell
further behind with \$127.35 for PPO, and \$74.68 for DQN.

However, it is important to note that the test period consists of
generally favourable market conditions. Therefore, the experiments
cannot simply be viewed as evidence that reinforcement learning failed,
as staying invested in generally rising markets is already a strong
strategy

Buy-and-hold is therefore a difficult benchmark to improve upon. A
learned strategy that reduces exposure must therefore focus on periods
when leaving the market is advantageous enough to compensate for the
return not gained while being out of the market.

The experiments provide limited evidence that the tested agents could
make this decision and overall did not demonstrate such behaviour.

\textbf{5.9.2 Additional State Features Did Not Lead to Better
Performance}

The state comparison did not show a consistent advantage from adding
drawdown to the nine-feature state. This is important because the state
representation was designed to include information about market
movements, portfolio position, wealth, and risk.

For REINFORCE and DQN, the changes in return were -\$4.12 and \$3.13,
respectively, in the nine-to-ten-feature comparison. The lack of
improvement in performance does not necessarily mean these features were
poorly selected. Instead, their usefulness depends on whether the
learning algorithm can use them to make better decisions.

This does not imply that the nine-feature state is universally
sufficient. It means that, in this particular environment and with this
wealth-change objective, additional market information produced no
measurable improvement.

The later risk-aware experiments explain the agents' change in behaviour
when the objective explicitly included drawdown, but this mainly led to
reduced market participation.

\textbf{5.9.3 The Algorithms Differed in Their Use of the State}

One of the clearest differences between the algorithms is behavioural
rather than financial.

REINFORCE and PPO often produced policies that were classified as
state-independent. Hence, their relatively strong returns need to be
carefully interpreted. A policy that buys whenever possible can perform
well during a period of rising prices without making significant use of
the market information provided to it.

Deep Q-learning showed a different pattern. It consistently produced
state-dependent behaviour in 6 of 6 seeds, indicating that its decisions
changed with the state currently being observed. However, this did not
lead to higher returns. Instead, the agent generally reduced its
exposure to the asset.

This result shows why financial return alone is not enough to evaluate
the learned trading policies from reinforcement learning algorithms.
Without the behavioural analysis, the strong performance of REINFORCE
could be mistaken for evidence that the model had learned a useful
relationship between market information and trading decisions.

\textbf{5.9.4 Lower Risk Was Mainly Associated with Lower Exposure}

The results show a consistent relationship between market participation
and risk; specifically, risk was reduced mainly by participating less in
the market.

Deep Q-learning achieved lower drawdown than buy-and-hold, but this was
due to lower exposure, and overall, it led to lower returns in the same
proportion.

This means lower drawdown should not be interpreted as evidence that DQN
avoided the worst market movements, as this differs from identifying and
avoiding specific periods of downside markets.

A proper risk-management strategy would ideally remain invested during
rising market periods while reducing exposure during market downsides
once detected. The current results do not provide enough evidence to
conclude that the agent learned this type of selective risk management.

\textbf{5.9.5 The Trade Size Affects Performance}

The trade size experiment shows that the transaction size is part of the
problem definition rather than merely an implementation design.

A small trading size limits how much exposure an agent can build during
a monthly episode. This creates a limitation when comparing the learned
agent with the simple buy-and-hold strategy. If the environment prevents
the agent from investing enough in the asset, poor performance may stem
from trade size rather than the quality of the learned policy.

Under the initial \(0.10C_{0}\) setting, the attainable exposure was
considerably lower than that of the buy-and-hold benchmark, reaching
only 0.690 average exposure as compared to 0.910 for buy-and-hold

The \(0.25C_{0}\) configuration addresses this limitation by increasing
the agent's attainable exposure to 0.838, giving a fairer basis for
comparing the learned policies with the benchmark, while retaining more
control over the positions held.

The trade size experiments also suggest that changing the trade size can
have a larger effect on performance than adding features to the state.
This supports the decision to calibrate trade size accurately before
working on the main learned strategies.

\textbf{5.9.6 Transaction Costs Provided a Useful Behavioural Test}

The transaction-fee experiment showed that Deep Q-learning reduced its
trading activity as transaction costs increased, suggesting that the
agent responds to the financial costs of its actions within the
environment.

Deep Q-learning\textquotesingle s trading activity reduced from 10.9
trades per month with no transaction fee to 1.9 trades at a 0.50\%
transaction fee, representing an approximately 83\% reduction.
REINFORCE, however, did not show the same response.

This provides further evidence that Deep Q-learning was making greater
use of the information available through the environment. However, this
additional responsiveness did not result in better financial
performance.

\textbf{5.9.7 Variation Between Seeds Is Important}

This experiment showed considerable variation between random seeds,
especially for the policy-gradient methods.

In some configurations, different seeds produced noticeably different
behaviours. Some policies turned towards high exposure, while others
traded less or became more conservative.

A seed producing a strong return would not be enough to conclude that
the algorithm had learned a reliable trading strategy. Consistent
behaviour across multiple seeds proves that, and becomes stronger
evidence that, the observed result is derived from that algorithm rather
than a product of one training run.

\textbf{5.9.7 Implications for the Research Objectives}

Overall, the experiments provide a qualified answer to the main research
question of whether reinforcement learning algorithms can learn
profitable trading strategies for a portfolio (in this case, the SPY
asset).

The reinforcement learning agents learned trading policies, but the
experiments did not show they could outperform the simpler buy-and-hold
benchmark on unseen data.

The results also show that learning behaviour strongly depends on the
algorithm, with clear differences between the learning methods.
REINFORCE moved towards remaining invested and made little use of state
information. DQN, however, consistently changed its actions based on the
observed state, making it highly state-dependent, but it did not yield
higher returns.

PPO behaviour was closer to REINFORCE than to DQN. Although PPO uses an
actor-critic structure and a clipped policy objective to stabilise
updates, its performance was less consistent across seeds, with a larger
spread in returns. This suggests that PPO\textquotesingle s additional
stability mechanisms did not produce an effective trading policy in this
environment.

The risk-aware reward also changed the behaviour of the state-dependent
agents and reduced drawdown, but the reduction was mainly due to lower
market exposure rather than better timing for entering and exiting
markets during downward periods. When the penalty became too large, the
agents almost stopped trading altogether.

These findings suggest that the main challenge identified by the
experiments is not just related to providing the agent with more
information. Instead, the results strongly point to the interaction
between the \textbf{reward function, trade size, and the behaviour
learned by the agent} as key factors in determining whether the
resulting strategy is useful.

\textbf{5.10 Chapter Summary}

The chapter first determines the final trade size because this affects
how quickly an agent can build or reduce its position. It then examines
whether the learned strategies rely on the information from the state
representations, followed by a comparison between the main nine-feature
state and the additional risk-aware state. It then presents the main
held-out test results, followed by experiments on the risk-aware reward
and the effect of transaction costs. The chapter concludes by bringing
the findings together and relating them to the research objectives.

The results show that none of the learned agents consistently
outperforms the buy-and-hold benchmark over the held-out dataset. The
behavioral analysis offers more insight into this outcome. The
policy-gradient algorithms rely less on available state information,
whereas the value-based agent uses more available state information but
generally has a lower exposure. The experiments also indicate that the
fixed trading amount affects performance. Finally, the drawdown-based
reward reduces both return and exposure, suggesting that the risk-aware
reward mainly encourages lower market participation rather than
generating better risk-adjusted trading performance.

\textbf{Chapter 6 -- Conclusion and Future Work.}

\textbf{6.1 Summary of the Work}

This dissertation examined whether reinforcement learning algorithms can
learn profitable trading policies that utilize market and portfolio
information from the observed states when trading assets in a portfolio.

The study designed a complete experimental system for trading the SPY
exchange-traded fund. The system uses daily market data organised into
monthly episodes, includes transaction costs and action masking, and
uses a nine-feature state representation with an optional drawdown
feature. Two reward formulations were tested: the primary wealth-change
reward and an optional risk-aware reward that penalised drawdown, along
with three learning algorithms: REINFORCE, Deep Q-learning (DQN), and
Proximal Policy Optimisation (PPO).

Experimental decisions were made on the 12-month validation dataset
(2023), while keeping the 24-month test set (2025-2025) separate and
using it only for final evaluations.

The experiments followed a fixed sequence. First, the trade size was
adjusted to ensure the agents could reach a reasonable level of exposure
relative to the buy-and-hold benchmark. The behavioural experiment then
determined whether the learned policies actually respond to the
information in their states. Once these environment, calibration and
behavioural choices had been decided, the main performance of the
learned agents was evaluated and compared against the benchmark strategy
on the unseen test period.

Additional experiments then tested whether adding drawdown to the state
changed the results, whether a risk-aware reward improved loss
management in the portfolio trade, and how the different learning
methods responded to changes in transaction costs.

\textbf{6.2 Conclusions Against the Research Questions}

The experiments in Chapter 5 address the four questions asked in 5.1.

\textbf{6.2.1 Performance}

The first question considered whether the reinforcement learning agents
could outperform the buy-and-hold benchmark on unseen data.

The results do not provide any evidence that they could. No learned
agent outperformed the simple buy-and-hold strategy on the held-out test
period. Buy-and-hold produced a mean wealth change of +\$180.25 per
month, while the best learned agent, REINFORCE, achieved +\$171.51. PPO
achieved +\$127.35 and DQN +\$74.68.

The result should not be interpreted as reinforcement learning
algorithms being incapable of learning profitable trading strategies.
All agents produced positive average wealth changes. Instead, the more
specific conclusion is that none of the tested algorithms learned a
level of trading skill that could outperform the simple investment
strategy on unseen data.

Transaction costs and the Trade size also cannot explain the difference,
since they were the same throughout the experiment, providing a basis
for fair comparison. The test period, however, favoured remaining
invested, with positive market performance in 17 of the 24 months. Under
these conditions, buy-and-hold was difficult to beat as it benefits from
the rising market direction.

\textbf{6.2.2 Behaviour}

The second question examined whether the learned policies genuinely use
the state information provided to them.

The algorithms showed a clear difference in how they used the state. All
six DQN seeds changed their actions in response to the observed state,
while none of the twelve policy-gradient seeds from REINFORCE and PPO
did so. This changes how the performance results should be interpreted.
REINFORCE came close to the benchmark in terms of return, at +\$171.51
per month, but its policies largely shifted to buying rather than using
the market information provided.

DQN however produced policies that consistently responded to changes in
the observed state but produced a lower average return than REINFORCE.
This suggests that state dependence and financial performance are not to
be treated as equal objectives in the experimental setting. An agent can
learn to use its observations without that behaviour necessarily
yielding higher returns.

This differentiation is important because an agent can perform well for
the wrong reason. During a generally rising market, a policy that
remains invested can achieve a strong return without needing to learn
when market conditions require entering, reducing, or exiting a
position.

Looking only at return would have completely hidden the distinction
between these behaviours. The behavioural analysis is therefore an
important part of evaluating learned trading policies, as it revealed a
clear difference between the value-based and policy-gradient methods.

\textbf{6.2.3 State Representation}

The third question examined whether adding a drawdown feature to the
state representation improved agent performance or behaviour.

Adding drawdown to the nine-feature state produced only limited changes
using the wealth-change reward. REINFORCE\textquotesingle s mean monthly
wealth change decreased by \$4.12, while DQN\textquotesingle s increased
by \$3.13. Both changes were small compared with the variation between
their random seeds and cannot be reasonably treated as reliable
improvements.

PPO\textquotesingle s mean return increased by \$32.62, but its seed
spread also changed heavily, from \$179.80 to \$94.00. This improvement
cannot confidently be attributed to the drawdown feature itself.

Overall, the results suggest that providing additional information does
not automatically produce a better policy or improve learning unless the
objective specifies it. This is consistent with an important principle
noticed from the experiments: \textbf{information in the state is only
useful when the learning objective and optimisation process give the
agent a reason to use it}. Hence, the agent must be able to learn a
relationship between the additional information and actions that improve
the objective being optimised.

\textbf{6.2.4 Reward Formulation}

The fourth question considered whether adding drawdown into the reward
could encourage the agents to reduce losses and generate a more
risk-aware behaviour.

The risk-aware reward reduced drawdown mainly by reducing market
participation, not by improving understanding of market conditions. With
the selected penalty of \(\mathbf{\lambda = 0.10}\),
DQN\textquotesingle s mean drawdown fell by 41\%, from 0.0140 to 0.0082.
However, its mean return also fell by 32\%, from +\$77.81 to +\$52.90,
while its exposure was approximately halved. Its Sharpe ratio fell from
0.505 to 0.416, so the reduced drawdown did not result in better
risk-adjusted performance. The reward therefore encouraged the agent to
hold less of the asset rather than to identify and avoid declining
markets. This is also reflected in March 2025, when reducing exposure
would have been particularly useful: DQN\textquotesingle s exposure
during the month was higher than its overall average, and it still
experienced 46\% of the benchmark\textquotesingle s loss.

The result illustrates an important limitation of using drawdown as a
standalone penalty. The easiest way for an agent to reduce exposure to
drawdown is to reduce its participation in the market. This behaviour
can theoretically satisfy the risk conditions of the reward without
necessarily producing a better ability to identify when the market is
likely to decline.

The penalty-weight analysis proves this interpretation. A relatively
small increase in the penalty produced a major behavioural change, with
exposure falling from 0.344 at λ = 0.10 to 0.074 at λ = 0.25. This
indicates that the reward formulation created a strong correlation
between market participation and risk avoidance.

\textbf{6.3 Critical Appraisal of the Project}

The most important lessons from this project were from the experimental
methodology rather than the choice of reinforcement learning algorithm
alone. Several of the most useful findings came from implementation
problems identified during development and testing.

One of the strongest aspects of the project was the verification tests.
The automated tests identified \textbf{12 implementation problems,}
three of which could have affected the experiments' validity. These
included issues with rolling feature calculation at the end of monthly
episodes and an implementation error in the buy-and-hold benchmark that
reduced its performance by \$38.95 per month, about 21.6\%. If this
error had remained, the dissertation\textquotesingle s main conclusion
could have been misleading, as the learned agents would have appeared to
outperform the benchmark.

The verification tests also noticed the experimental design problem in
the trade size. Using the original \(0.10C_{0}\) trade size in Chapter 3
would make the maximum attainable exposure only 0.690, compared with
0.910 for buy-and-hold. This meant that the learning agents were at a
disadvantage before training even began. Selecting \(0.25C_{0}\) allowed
a fairer comparison while still letting the agents control their
positions in smaller steps than the\(0.50C_{0}\) alternative.

The complete set of 21 verification tests passed before the final
experiments were accepted, giving a higher confidence that the reported
results reflect an appropriate experimental design rather than
implementation errors.

The behavioural analysis was another important part of the study.
Without it, the REINFORCE result of +\$171.51 could have been
interpreted as evidence that the agent had learned an intelligent,
effective trading strategy. The analysis showed this interpretation was
not justified, as the policies acted like a consistent buyer because the
market was consistently rising, completely changing the interpretation
of the result.

The project also highlighted the importance of random seed variation.
The results, particularly for PPO, varied considerably between
independent training runs. PPO had a seed spread of \$179.80 in one
configuration, which was greater than its mean return. This made it
difficult to draw conclusions based on individual training runs, as that
would have been unreliable. Reporting multiple seeds widened the
evaluation, providing a more consistent view of how each algorithm
behaved, although it increased the computational and experimental
workload; however, six seeds are still lacking for more definite
statistical inference.

Overall, the project demonstrated that careful experimental design is as
important as the choice of reinforcement learning algorithm. The main
lessons came from testing the assumptions behind the experiments, rather
than from simply comparing model returns. State dependence, action
constraints, reward design, and seed variation all affect how the final
results could be interpreted. This made the experimental process more
demanding but also produced a more reliable basis for evaluating the
learned trading policies.

\textbf{6.4 Limitations}

The findings of this study should be considered within the scope of the
experimental design, but there are several limitations.

\textbf{6.4.1 One asset and a limited market period}

The experiments use SPY as the only asset, and the held-out period from
2024 to 2025 was generally favourable to holding the asset due to rising
markets. A market with a longer or more severe decline could produce
different results for the learned agents and the buy-and-hold benchmark.
The findings cannot therefore be automatically generalised to other
assets or market conditions without further testing.

\textbf{6.4.2 Fixed Trade Size}

The agents trade fixed fractions of the initial capital rather than
choosing their position size continuously. This was calibrated to create
a controlled, reasonable environment for comparison with the benchmark,
but it restricts how much the agents can express their decisions. A
continuous position-sizing model could potentially respond differently
based on its exposure in a particular market condition. Rather than
repeatedly applying fixed buy or sell actions, such an agent could
select its desired portfolio exposure, producing different results in
returns and trading behaviour.

\textbf{6.4.3 Limited Number of Seeds and Lack of Significance Testing}

Six seeds were used for each configuration, and the seed-level results
and variation between them are reported throughout the results. This is
stronger than relying on a single training run, but it does not provide
sufficient statistics for formal claims about small differences. For
example, the \$3.13 change observed for DQN and \textbf{−\$4.12} for
REINFORCE when the state representation was changed is not large enough
to be treated as a meaningful improvement.

\textbf{6.4.5 Risk Penalty Selection was based on DQN}

The reward\textquotesingle s penalty weight λ was selected using DQN
because it was the algorithm that consistently showed state-dependent
behaviour. This does not guarantee that the same penalty value is
appropriate for REINFORCE or PPO, as it was not independently optimised
for every algorithm. A separate calibration for each algorithm could
produce different results. This limitation is particularly relevant
because the algorithms responded differently to the state and reward
formulation.

\textbf{6.4.6 The risk penalty only targets drawdown.}

The risk-aware reward penalises drawdown but does not directly measure
whether the return achieved is sufficient for the level of risk taken.
While drawdown is directly relevant to losses, it does not fully
represent the portfolio risk. A reward that relates return to risk,
rather than penalizing risk alone, was not tested.

Additionally, as shown in Section 5.7, the agents could reduce the
penalty simply by reducing their exposure to the asset or not trading at
all. The experiments showed that a lower drawdown value should not
automatically be associated with evidence of better risk management.

\textbf{6.5 Future Work}

The limitations above suggest several directions for furthering the
research. Each stem directly from a limitation or an unresolved result
in the experiments. The behavioural analysis should also be applied to
these future works to determine the agent's behaviour to changes
alongside performance improvements.

\textbf{6.5.1 Longer and More Challenging Market Periods}

The experiments could be repeated over longer market periods containing
more market declines. A useful change would be to use a walk-forward
training approach, allowing the agents to be retrained as market
conditions change. The current study provides useful evidence for a
generally rising test period, but testing market declines would provide
a balanced evaluation of the portfolio.

This would better reflect the changing nature of financial markets,
allowing the system to encounter different combinations of rising,
falling, and volatile conditions. Such an evaluation would also better
test the agents' risk-aware behaviour during downtrends.

\textbf{6.5.2 Predictive and Uncertainty-Aware State Information}

The drawdown feature describes the portfolio\textquotesingle s previous
losses but does not provide information about future market movements.
Future work could therefore include predictive information, such as
probabilistic forecasts containing estimates of forecast uncertainty.

This would give a state-dependent agent more useful information for
making decisions about future market conditions rather than relying only
on historical data. It would also directly test whether
DQN\textquotesingle s ability to use the state leads to more useful
market understanding and trading decisions.

\textbf{6.5.3 Risk-Efficiency Objectives}

The most direct correction to the drawdown penalty could be to replace
it with a reward that considers return relative to risk, such as a
Sharpe-based or drawdown-relative objective. This would make simply
reducing market exposure a less effective way of improving the reward,
as seen in the current experiments

The results in Section 5.7 make this future rework highly relevant
because the tested penalty mainly caused agents to reduce exposure
rather than learn when downside risk is increasing.

\textbf{6.5.5 More Seeds for Better Statistical Analysis}

More random training seeds would improve the reliability of comparisons
between algorithms and configurations\textbf{.} Future work should
combine this with appropriate statistical analysis to determine whether
observed differences between configurations are likely to represent
genuine results rather than random noise.

\textbf{6.5.6 More flexible trade sizes}

Future experiments could allow the agent to choose its position size
continuously rather than using fixed trade amounts. This would remove
the exposure limit created by the current environment and allow the
agent to express different levels of confidence by adjusting its
position size, potentially enabling smoother portfolio adjustments.

\textbf{6.6 Closing Remarks}

This dissertation examined whether standard reinforcement learning could
beat the simple buy-and-hold strategy on unseen data. The experiments
found no evidence that it could. The best-performing learned agent
finished \$8.74 per month below buy-and-hold, and it did so by ignoring
their market observations to follow a simple trading behaviour.

The main contribution of the study is therefore not a successful trading
strategy, but a careful examination of why the learned agents performed
as they did. The comparisons were conducted in a controlled environment
in which trade sizing was calibrated, transaction costs were applied,
and the final results were produced in a suitable experimental
framework. The policies were then evaluated on both behavioral and
financial performance, and the main findings were tested across the
different state representations and reward functions.

The implementation, experimental framework and behavioural analysis
developed in this study provide a useful basis for future work. More
advanced models can build on this framework while being evaluated and
compared against the same fair standards.
